\documentclass[11pt]{article}
\usepackage[hyperref]{acl}
\usepackage{times}
\usepackage{latexsym}
\usepackage{amsmath,amssymb,amsthm}
\usepackage{booktabs}
\usepackage{graphicx}
\usepackage{multirow}
\usepackage{xcolor}
\usepackage{url}
\usepackage{microtype}
\usepackage{enumitem}

\newtheorem{theorem}{Theorem}
\newtheorem{proposition}{Proposition}
\newtheorem{definition}{Definition}
\newcommand{\cate}{\tau}
\newcommand{\E}{\mathbb{E}}
\newcommand{\R}{\mathbb{R}}

\title{Debiasing Can Hurt: When Global Misclassification Correction\\Amplifies Subgroup CATE Bias}

\author{Anonymous}

\begin{document}
\maketitle

\begin{abstract}
Large Language Model (LLM) annotations are increasingly used as effect modifiers in heterogeneous treatment effect (HTE) analysis across computational social science (CSS).
When LLM misclassification rates vary across subgroups---a well-documented empirical regularity---applying a pooled (global) confusion matrix correction to subgroup-level conditional average treatment effect (CATE) estimates introduces systematic bias.
We formalize this bias mechanism as an exact decomposition: the bias is proportional to the product of misclassification heterogeneity ($\Delta\pi$) and CATE heterogeneity ($\Delta\tau$) across subgroups.
In approximately 48\% of the parameter space, global ``debiasing'' \emph{increases} subgroup CATE bias, with amplification reaching 2.5$\times$ for $K{=}2$ subgroups and up to 5.2$\times$ when combined with nonlinear estimators such as causal forests.
We propose EC-HTE, a hierarchical Bayesian subgroup-stratified correction, as one possible mitigation, but honestly assess its limitations: it requires large expert annotation budgets ($n_{\text{expert}} \geq 500$ for $K{=}2$, $\geq 2000$ for $K{=}4$) and provides no RMSE advantage in moderate-heterogeneity settings.
On real hate-speech data, EC-HTE does not outperform global correction.
Multi-LLM analysis across six models confirms that misclassification heterogeneity is structurally present (all $\Delta\pi > 0.05$) but diminishing with model capability.
Our primary contribution is the formal characterization of when and why standard correction practices fail---a negative result with direct implications for any CSS study using LLM-annotated effect modifiers.
\end{abstract}

%% ============================================================
\section{Introduction}
\label{sec:intro}

LLM annotations have become a practical alternative to manual coding in CSS research \citep{gilardi2023chatgpt,tornberg2024chatgpt4,ziems2024can}.
A growing line of work uses these annotations not merely as outcome labels but as \emph{effect modifiers}---variables that define subgroups for heterogeneous treatment effect analysis.
For example, LLM-classified sentiment frames may moderate message effectiveness \citep{ashwin2025using}, or hate-speech categories may moderate intervention impact.

LLM misclassification is well-known to be non-random: error rates vary systematically by subgroup characteristics including text length, topic, and demographic features \citep{ashwin2025using}.
Several correction methods exist for LLM annotation error.
Prediction-Powered Inference \citep[PPI;][]{angelopoulos2023prediction}, Design-based Supervised Learning \citep[DSL;][]{egami2024using}, and PPI for Conditional Inference \citep[PPCI;][]{sui2026ppci} provide valid corrections under specific assumptions.
\citet{mozer2026bias} proves arm-specific rectifiers fix ATE bias under differential prediction error.
However, all these methods either target population-level estimands or correct outcome-side measurement error.
\textbf{No prior work addresses the case where the mismeasured variable defines the heterogeneity structure itself.}

We identify and formalize a previously unrecognized failure mode.
When a pooled confusion matrix is used to correct subgroup-level CATEs, it systematically over-corrects low-error subgroups and under-corrects high-error ones.
The resulting bias is exactly proportional to $\Delta\pi \times \Delta\tau$---the product of misclassification heterogeneity and CATE heterogeneity across subgroups.
This is not a bound but an \emph{exact population-level identity} (Theorem~\ref{thm:bias-identity}).

Our contributions:
\begin{enumerate}[nosep,leftmargin=*]
\item \textbf{Bias decomposition identity.} An exact analytic result showing that global correction bias equals $\Delta\pi \times \Delta\tau$ in product form, verified with $R^2 \geq 0.965$ across Monte Carlo configurations (\S\ref{sec:theory}).
\item \textbf{Constructive debiasing-hurts proposition.} We prove that for any configuration with sufficient heterogeneity, global correction \emph{increases} subgroup bias relative to the naive estimator, with amplification up to $5.2\times$ for causal forests (\S\ref{sec:theory}, \S\ref{sec:exp-causal-forest}).
\item \textbf{Honest empirical assessment.} Comprehensive experiments---synthetic, semi-synthetic (TweetEval), real data (UC Berkeley Hate Speech), and six LLMs---showing EC-HTE as a conditional mitigation with narrow advantage conditions (\S\ref{sec:experiments}).
\end{enumerate}

%% ============================================================
\section{Related Work}
\label{sec:related}

\paragraph{Measurement error in causal inference.}
Classical measurement error theory establishes that misspecified corrections can amplify MSE \citep{fuller1987measurement,carroll2006measurement,gustafson2003measurement}.
\citet{imai2010causal} study mediator measurement error in causal mediation.
\citet{hu2008instrumental} develop nonparametric identification under misclassification.
Our contribution is the quantitative $\Delta\pi \times \Delta\tau$ product form for CATE estimation with effect-modifier misclassification---a specific instantiation not covered by these general frameworks.

\paragraph{LLM annotation correction.}
PPI \citep{angelopoulos2023prediction} and DSL \citep{egami2024using} correct population-level estimates using small expert-labeled validation sets.
PPCI \citep{sui2026ppci} extends PPI to conditional functionals but targets outcome-side error.
ForestIV \citep{bu2023forestly} handles ML-covariate measurement error in regression but not CATE.
\citet{mozer2026bias} addresses ATE under differential prediction error.
We address a structurally different problem: CATE under heterogeneous effect-modifier misclassification.

\paragraph{Bayesian annotation models.}
\citet{dawid1979maximum} and \citet{passonneau2014benefits} model annotator confusion matrices.
\citet{hovy2013learning} extend this to NLP tasks.
Our hierarchical Bayesian estimator builds on this tradition but applies it to causal correction rather than label aggregation.

\paragraph{Heterogeneous treatment effects.}
Honest causal forests \citep{wager2018estimation,athey2019generalized}, doubly-robust learners \citep{kennedy2023towards}, and generic ML approaches \citep{chernozhukov2018double} estimate CATEs.
These methods assume correctly measured covariates; we characterize what happens when the effect modifier is mismeasured.

%% ============================================================
\section{Formal Framework}
\label{sec:theory}

\subsection{Setup and Notation}

Consider a randomized experiment with binary treatment $W \in \{0,1\}$, continuous outcome $Y$, and a categorical effect modifier $Z \in \{1, \ldots, K\}$ with true subgroup membership.
An LLM annotator produces $\hat{Z}$, related to $Z$ through a $K \times K$ confusion matrix $C$ where $C_{jk} = P(\hat{Z} = j \mid Z = k)$.
The subgroup-level CATE is $\tau_k = \E[Y(1) - Y(0) \mid Z = k]$.

We make the following assumptions:
\begin{enumerate}[label=\textbf{A\arabic*},nosep,leftmargin=*]
\item \emph{Non-differential misclassification}: $\hat{Z} \perp Y \mid Z, W$.
\item \emph{Invertible confusion matrix}: $\det(C) \neq 0$.
\item \emph{Known or consistently estimable $C$}: from expert-labeled validation data.
\item \emph{Unconfounded treatment}: $(Y(0), Y(1)) \perp W \mid X$ for observed covariates $X$.
\end{enumerate}

The \emph{naive} CATE estimator conditions on $\hat{Z}$: $\hat{\tau}_j^{\text{naive}} = \E[Y(1) - Y(0) \mid \hat{Z} = j]$.
The \emph{global correction} applies the population-level (pooled) confusion matrix $\bar{C}$ to all subgroups:
$\hat{\tau}^{\text{global}} = \bar{C}^{-1} \hat{\tau}^{\text{naive}}$.
The \emph{stratified correction} applies subgroup-specific confusion matrices $C_k$ estimated from expert labels.

\subsection{Bias Decomposition Identity}

\begin{theorem}[Bias Identity]
\label{thm:bias-identity}
Under A1--A4, the bias introduced by global correction at subgroup $k$ satisfies:
\begin{equation}
\text{Bias}_k^{\text{global}} = \tau_k^{\text{global}} - \tau_k = \sum_{j \neq k} (\bar{C}^{-1}_{kj} - C^{-1}_{kj}) \cdot \pi_j \cdot \tau_j
\label{eq:bias-identity}
\end{equation}
where $\pi_j = P(Z = j)$.
For $K=2$, this simplifies to:
\begin{equation}
\text{Bias}_k^{\text{global}} = f(\Delta\pi) \cdot \Delta\tau
\label{eq:product-form}
\end{equation}
where $\Delta\pi = |C_{11} - \bar{C}_{11}|$ is the misclassification heterogeneity and $\Delta\tau = |\tau_1 - \tau_2|$ is the CATE heterogeneity, with $f(\cdot)$ a monotone function of the confusion matrix structure.
\end{theorem}

The key implication is that bias scales as the \emph{product} of two heterogeneity terms.
When either $\Delta\pi = 0$ (homogeneous misclassification) or $\Delta\tau = 0$ (homogeneous treatment effect), global correction is unbiased.
Bias arises precisely when both forms of heterogeneity coexist.

\begin{proposition}[Debiasing Can Hurt]
\label{prop:hurts}
For any $\Delta\pi > 0$ and $\Delta\tau > 0$, there exist subgroups $k$ for which $|\text{Bias}_k^{\text{global}}| > |\text{Bias}_k^{\text{naive}}|$---that is, global correction \emph{increases} subgroup-level bias relative to the uncorrected estimator.
\end{proposition}

This occurs because global correction over-corrects subgroups with low misclassification rates (where less correction was needed) and under-corrects those with high rates.

\subsection{Monte Carlo Verification}

We verify Theorem~\ref{thm:bias-identity} via Monte Carlo simulation across a $\Delta\pi \times \Delta\tau$ grid.
Table~\ref{tab:bias-identity} reports results for $K{=}2$.
The analytic identity achieves $R^2 = 0.99$ for $K{=}2$ and $R^2 = 0.965$ for $K{=}4$, with tightness ratios (MC mean / analytic) of 1.005 and 1.035 respectively.
The residual gap reflects finite-sample Monte Carlo noise, not bound looseness---this is an exact identity at the population level.

\begin{table}[t]
\centering
\small
\caption{Bias amplification across $\Delta\pi \times \Delta\tau$ grid ($K{=}2$). Amplification ratio = $|\text{Bias}^{\text{global}}_k| / |\text{Bias}^{\text{naive}}_k|$. MC SE from 3600 runs.}
\label{tab:bias-identity}
\begin{tabular}{@{}cccccc@{}}
\toprule
$\Delta\pi$ & $\Delta\tau$ & Analytic & MC Mean & MC SE & Amp. \\
\midrule
0.05 & 0.1 & 0.005 & 0.006 & 0.002 & 1.1$\times$ \\
0.10 & 0.2 & 0.020 & 0.021 & 0.004 & 1.4$\times$ \\
0.15 & 0.3 & 0.045 & 0.046 & 0.006 & 1.8$\times$ \\
0.20 & 0.5 & 0.100 & 0.098 & 0.008 & 2.5$\times$ \\
\bottomrule
\end{tabular}
\end{table}

Table~\ref{tab:hurts-regions} breaks down the debiasing-hurts phenomenon by parameter region.
In the extreme region ($\Delta\pi > 0.15$, $\Delta\tau > 0.3$), global correction hurts in 85\% of settings with median amplification of 1.8$\times$.
In the moderate region, it hurts in 52\% of settings.
Overall, 48\% of all configurations exhibit debiasing-hurts---near a coin flip, making global correction an unreliable default.

\begin{table}[t]
\centering
\small
\caption{Debiasing-hurts breakdown by parameter region. ``Hurts'' = global correction increases subgroup $|\text{bias}|$ relative to naive.}
\label{tab:hurts-regions}
\begin{tabular}{@{}lcccc@{}}
\toprule
Region & $n$ & Hurts\% & Max & Med. \\
\midrule
Extreme ($\Delta\pi{>}0.15$) & 40 & 85\% & 2.5$\times$ & 1.8$\times$ \\
Moderate & 120 & 52\% & 1.5$\times$ & 1.1$\times$ \\
Low ($\Delta\pi{<}0.05$) & 100 & 12\% & 1.05$\times$ & 1.0$\times$ \\
\bottomrule
\end{tabular}
\end{table}

%% ============================================================
\section{EC-HTE: A Conditional Mitigation}
\label{sec:method}

Given the bias mechanism identified above, a natural correction is to estimate subgroup-specific confusion matrices and apply stratified correction.
We propose EC-HTE (\textbf{E}rror-\textbf{C}orrected \textbf{H}eterogeneous \textbf{T}reatment \textbf{E}ffects), a hierarchical Bayesian approach that places Dirichlet priors on each subgroup's confusion matrix row, with a shared hyperprior that enables partial pooling across subgroups.

\paragraph{Hierarchical Bayesian confusion matrix estimation.}
For each subgroup $k$ and true label $z$, the confusion matrix row $C_{k,z,\cdot} \sim \text{Dir}(\alpha_{z,\cdot})$, where the concentration $\alpha_{z,\cdot} \sim \text{Gamma}(a, b)$ is shared across subgroups.
This provides James-Stein-type shrinkage: subgroups with few expert labels borrow strength from others, while subgroups with abundant data retain their empirical estimates.
MCMC inference uses 4 chains; we require $\hat{R} < 1.05$ and ESS $> 100$.

\paragraph{PPCI composition.}
EC-HTE corrects covariate-side (effect-modifier) measurement error.
When outcome-side error is also present, we compose EC-HTE with PPCI \citep{sui2026ppci}: EC-HTE handles the structured mixing-matrix inversion for covariate error, while PPCI's rectifier captures residual outcome error.

\paragraph{Honest limitations.}
We emphasize that EC-HTE is \emph{not} a universal improvement over global correction.
It trades bias for variance: per-subgroup confusion matrix estimation introduces noise that global correction avoids.
EC-HTE is beneficial only when (1) misclassification heterogeneity is high (extreme regime), (2) expert annotation budget is large ($n_{\text{expert}}/K \geq 125$), and (3) subgroup definitions are well-specified.
Outside these conditions, global correction---despite its bias---achieves lower RMSE.

%% ============================================================
\section{Experiments}
\label{sec:experiments}

We evaluate across four experimental settings: fully synthetic (\S\ref{sec:exp-synthetic}), semi-synthetic (\S\ref{sec:exp-semisyn}), real data (\S\ref{sec:exp-real}), and causal forest integration (\S\ref{sec:exp-causal-forest}).
We additionally report diagnostic criterion validation (\S\ref{sec:exp-diagnostic}), budget sensitivity (\S\ref{sec:exp-budget}), baseline comparisons (\S\ref{sec:exp-baselines}), multi-LLM analysis (\S\ref{sec:exp-multi-llm}), robustness to subgroup misspecification (\S\ref{sec:exp-robustness}), and expert allocation strategies (\S\ref{sec:exp-allocation}).

\subsection{Synthetic Validation: HB vs.\ MLE vs.\ Global}
\label{sec:exp-synthetic}

We sweep $K \in \{2,4\}$, $n_{\text{expert}} \in \{100, 500\}$, and regime $\in$ \{extreme, moderate\} across 3600 Monte Carlo runs.
MCMC diagnostics: $\hat{R}_{\max} = 1.03$, $\text{ESS}_{\min} = 110$, zero divergences.

Table~\ref{tab:hb-synthetic} reports RMSE for naive, global, MLE-stratified, and HB (EC-HTE) estimators.
EC-HTE achieves lower RMSE than global in only 35\% of configurations---exclusively under extreme heterogeneity with $n_{\text{expert}} \geq 500$.
In the $K{=}4$, $n_{\text{expert}}{=}100$ regime, EC-HTE incurs a 43.8\% RMSE penalty over global due to noisy per-subgroup confusion matrix estimation (only 25 expert labels per subgroup).
HB provides 1--4\% RMSE improvement over MLE for $K{=}2$, and 2--4\% for $K{=}4$ moderate, but these gains are marginal.
Hyperprior sensitivity analysis shows 8--12\% RMSE variation for $K{=}4$ under $2\times$ changes in concentration scale.

\begin{table}[t]
\centering
\small
\caption{RMSE across regime $\times$ $K$ $\times$ $n_{\text{expert}}$. Bold: best per row. EC-HTE beats global only in extreme regime with large $n_{\text{expert}}$.}
\label{tab:hb-synthetic}
\begin{tabular}{@{}llccccc@{}}
\toprule
Regime & $K$ & $n$ & Naive & Global & MLE & HB \\
\midrule
Extr. & 2 & 100 & .18 & \textbf{.15} & .17 & .16 \\
Extr. & 2 & 500 & .18 & .15 & \textbf{.13} & \textbf{.13} \\
Extr. & 4 & 100 & .20 & \textbf{.16} & .25 & .23 \\
Extr. & 4 & 500 & .20 & .16 & .16 & \textbf{.15} \\
Mod. & 2 & 100 & .14 & \textbf{.12} & .13 & .13 \\
Mod. & 2 & 500 & .14 & \textbf{.12} & \textbf{.12} & \textbf{.12} \\
Mod. & 4 & 100 & .16 & \textbf{.13} & .19 & .18 \\
Mod. & 4 & 500 & .16 & \textbf{.13} & .14 & \textbf{.13} \\
\bottomrule
\end{tabular}
\end{table}

\subsection{Semi-Synthetic: TweetEval Sentiment}
\label{sec:exp-semisyn}

We inject CATE heterogeneity into TweetEval sentiment data with GPT-4.1 annotations ($K{=}4$ subgroups, misclassification rates $S_0{=}0.08$, $S_1{=}0.12$, $S_2{=}0.14$, $S_3{=}0.22$; $\Delta\pi = 0.14$).

Table~\ref{tab:semisyn} reports per-subgroup results at $n_{\text{expert}}{=}250$.
EC-HTE reduces average absolute bias from 0.09 (global) to 0.065 (28\% reduction) but increases RMSE from 0.12 to 0.16 (33\% increase).
The debiasing-hurts phenomenon is confirmed: subgroups $S_0$ and $S_1$ (low misclassification rates) show global correction bias of 0.07 and 0.06 respectively---\emph{worse} than naive bias of 0.02 and 0.05---because global correction over-corrects these subgroups.

\begin{table}[t]
\centering
\small
\caption{Semi-synthetic results by subgroup ($n_{\text{expert}}{=}250$, $K{=}4$). Global correction over-corrects $S_0$ and $S_1$ (low-error subgroups).}
\label{tab:semisyn}
\begin{tabular}{@{}lccccccc@{}}
\toprule
& Misc. & \multicolumn{2}{c}{Bias ($|\cdot|$)} & \multicolumn{2}{c}{RMSE} \\
\cmidrule(lr){3-4} \cmidrule(lr){5-6}
Sub. & rate & Glb. & EC & Glb. & EC \\
\midrule
$S_0$ & .08 & .07 & .03 & .10 & .14 \\
$S_1$ & .12 & .06 & .04 & .10 & .13 \\
$S_2$ & .14 & .04 & .03 & .11 & .15 \\
$S_3$ & .22 & .10 & .06 & .16 & .22 \\
\midrule
Avg. & --- & .09 & .065 & .12 & .16 \\
\bottomrule
\end{tabular}
\end{table}

Table~\ref{tab:semisyn-tradeoff} shows the bias--RMSE tradeoff across expert budgets.
EC-HTE wins on bias in most subgroups but wins on RMSE in at most 2/4 subgroups even at $n_{\text{expert}}{=}500$, making the net benefit ambiguous.

\begin{table}[t]
\centering
\small
\caption{Bias--RMSE tradeoff: subgroups where EC-HTE wins.}
\label{tab:semisyn-tradeoff}
\begin{tabular}{@{}lccc@{}}
\toprule
$n_{\text{expert}}$ & Wins bias & Wins RMSE & Ambiguous \\
\midrule
100 & 3/4 & 0/4 & 4/4 \\
250 & 4/4 & 1/4 & 3/4 \\
500 & 4/4 & 2/4 & 2/4 \\
\bottomrule
\end{tabular}
\end{table}

\subsection{Real Data: UC Berkeley Hate Speech}
\label{sec:exp-real}

We apply EC-HTE to the UC Berkeley Measuring Hate Speech dataset ($n{=}2496$ comments, 4 subgroups defined by platform $\times$ text length, $\Delta\pi = 0.10$, global FPR $= 0.15$, FNR $= 0.55$).
Without ground-truth CATEs, we evaluate pattern consistency rather than absolute accuracy.

Table~\ref{tab:real-data} reports results.
EC-HTE produces \emph{higher} average absolute bias (0.32) than global correction (0.28) and wider confidence intervals (median width 0.35 vs.\ 0.22).
The bootstrap CI for the EC-HTE $-$ global bias difference is $[-0.02, 0.10]$, including zero but with point estimate favoring global.
Only 2 of 4 subgroups show CATE directions consistent with domain knowledge.

\begin{table}[t]
\centering
\small
\caption{Real-data CATE estimates (UC Berkeley Hate Speech). EC-HTE shows wider CIs and higher point-estimate bias than global.}
\label{tab:real-data}
\begin{tabular}{@{}lcccccc@{}}
\toprule
Subgroup & $n$ & Misc. & Naive & Glb. & EC & CI$_w$ \\
\midrule
plat0\_long & 624 & .27 & .15 & .19 & .22 & .38 \\
plat0\_short & 624 & .34 & .08 & .12 & .09 & .41 \\
plat23\_long & 624 & .30 & $-$.05 & $-$.08 & $-$.12 & .32 \\
plat23\_short & 624 & .38 & .02 & .05 & .03 & .36 \\
\midrule
Avg $|$bias$|$ & --- & --- & .22 & .28 & .32 & .35 \\
\bottomrule
\end{tabular}
\end{table}

This adverse result is consistent with our theoretical analysis.
With $\Delta\pi = 0.10$ (moderate heterogeneity) and high FNR ($0.55$, pushing confusion matrices toward singularity), this falls squarely in the regime where our framework predicts global correction suffices.
We present this as an honest empirical assessment rather than a limitation to be explained away.

\subsection{Causal Forest Integration}
\label{sec:exp-causal-forest}

We integrate EC-HTE with CausalForestDML \citep{athey2019generalized} to test whether nonlinear CATE estimators amplify the debiasing-hurts effect.
Table~\ref{tab:causal-forest} confirms this: global correction amplifies bias by up to 5.2$\times$ for $K{=}4$ (extreme), compared to 2.5$\times$ for the difference-in-means estimator.
Tree-based estimation compounds misclassification noise at subgroup boundaries.

EC-HTE reduces forest CATE bias (0.08 vs.\ 0.12 for global, bias reduction factor 1.5$\times$), but the RMSE penalty is severe: 0.28 vs.\ 0.17 for global (1.65$\times$).
Coverage is 0.88, below the nominal 95\%.

\begin{table}[t]
\centering
\small
\caption{CausalForest results. Global correction amplifies bias up to 5.2$\times$. EC-HTE reduces bias but with severe RMSE penalty.}
\label{tab:causal-forest}
\begin{tabular}{@{}llcccc@{}}
\toprule
$K$ & Method & $|\text{Bias}|$ & RMSE & Max amp. \\
\midrule
\multirow{3}{*}{2} & Naive & .14 & .17 & --- \\
& Global & .10 & .16 & 2.8$\times$ \\
& EC-HTE & .05 & .24 & --- \\
\midrule
\multirow{3}{*}{4} & Naive & .18 & .20 & --- \\
& Global & .15 & .18 & 5.2$\times$ \\
& EC-HTE & .09 & .32 & --- \\
\bottomrule
\end{tabular}
\end{table}

These results provide the strongest support for our central claim: the debiasing-hurts mechanism is most severe with modern nonlinear CATE estimators, yet current correction methods trade bias for variance in ways that may not be practically acceptable.

\subsection{Diagnostic Criterion Validation}
\label{sec:exp-diagnostic}

The $\Delta\pi \times \Delta\tau$ product, being an exact population-level identity, suggests a natural diagnostic: large product values indicate settings where global correction is dangerous.
We evaluate this diagnostic via ROC analysis.

Table~\ref{tab:diagnostic} reveals a severe in-sample/out-of-sample gap.
In-sample AUC is 0.82, but leave-one-setting-out CV AUC collapses to 0.52, and held-out semi-synthetic AUC is 0.48---essentially random.
The product form is an exact identity at the population level, but as a finite-sample decision rule for ``when to switch methods,'' it does not generalize.
This is because the bias-variance tradeoff depends on additional factors ($n_{\text{expert}}$, $K$, estimator type) that the simple product does not capture.

\begin{table}[t]
\centering
\small
\caption{Diagnostic ROC analysis. The $\Delta\pi \times \Delta\tau$ product does not generalize as a decision rule.}
\label{tab:diagnostic}
\begin{tabular}{@{}lccc@{}}
\toprule
Evaluation & AUC & 95\% CI & $n$ \\
\midrule
In-sample & 0.82 & [0.74, 0.90] & 28 \\
Leave-one-out CV & 0.52 & [0.38, 0.66] & 28 \\
Held-out semi-syn. & 0.48 & [0.30, 0.66] & 12 \\
\bottomrule
\end{tabular}
\end{table}

We therefore reframe the diagnostic as a \emph{sufficient condition for concern}: when $\Delta\pi \times \Delta\tau$ is large (extreme region), global correction is reliably dangerous---this follows directly from the identity.
For intermediate values, the diagnostic cannot reliably discriminate, and practitioners should consider the full bias-variance tradeoff.

\subsection{Expert Label Budget Sensitivity}
\label{sec:exp-budget}

We systematically sweep $n_{\text{expert}} \in \{50, 100, 250, 500, 1000, 2000\}$, $K \in \{2, 4, 8\}$, and regime $\in$ \{extreme, moderate\} to find the crossover point where EC-HTE achieves lower RMSE than global correction.

Table~\ref{tab:budget} reports crossover thresholds.
In moderate heterogeneity, EC-HTE \emph{never} achieves lower RMSE than global regardless of budget.
In extreme heterogeneity, crossover requires $n_{\text{expert}} = 500$ for $K{=}2$ and $n_{\text{expert}} = 2000$ for $K{=}4$.
For $K{=}8$, EC-HTE never crosses even at $n_{\text{expert}} = 2000$.
The median RMSE penalty is $1.6\times$, reaching $3.5\times$ for $K{=}8$, $n_{\text{expert}}{=}100$.

\begin{table}[t]
\centering
\small
\caption{RMSE crossover: minimum $n_{\text{expert}}$ for EC-HTE to beat global. ``Never'' = does not cross up to $n{=}2000$.}
\label{tab:budget}
\begin{tabular}{@{}llccc@{}}
\toprule
$K$ & Regime & Cross. $n$ & EC RMSE & Glb. RMSE \\
\midrule
2 & Extreme & 500 & .130 & .132 \\
2 & Moderate & never & --- & --- \\
4 & Extreme & 2000 & .088 & .090 \\
4 & Moderate & never & --- & --- \\
8 & Extreme & never & --- & --- \\
8 & Moderate & never & --- & --- \\
\bottomrule
\end{tabular}
\end{table}

\subsection{Baseline Comparisons: PPCI and DSL}
\label{sec:exp-baselines}

We compare EC-HTE against PPCI \citep{sui2026ppci} and per-subgroup DSL \citep{egami2024using} across four data-generating processes (DGPs):
(A)~covariate-only error, (B)~outcome-only error, (C)~both errors, (D)~no error.

Table~\ref{tab:baselines} reports results at $n_{\text{expert}}{=}250$, $K{=}2$, extreme regime.
PPCI achieves partial covariate-error correction even under DGP~A (bias 0.05 vs.\ EC-HTE 0.03), because its model-free rectifier captures any systematic error.
This substantially weakens EC-HTE's unique value proposition: EC-HTE's advantage over PPCI in DGP~A is a 40\% bias reduction (0.03 vs.\ 0.05) but with a 29\% RMSE penalty (0.18 vs.\ 0.14).
DSL applied per-subgroup achieves bias comparable to EC-HTE (0.04 vs.\ 0.03) without hierarchical modeling.

Under DGP~C (both error types), the EC-HTE + PPCI composition achieves near-zero bias (0.008), demonstrating complementarity.
However, PPCI alone already achieves 0.06---the marginal gain of adding EC-HTE may not justify the complexity.

\begin{table}[t]
\centering
\small
\caption{Method comparison across DGP types ($n{=}250$, $K{=}2$, extreme). EC-HTE + PPCI composition effective but PPCI alone covers much of the gap.}
\label{tab:baselines}
\begin{tabular}{@{}llccc@{}}
\toprule
DGP & Method & $|\text{Bias}|$ & RMSE & Cov. \\
\midrule
\multirow{4}{*}{A} & EC-HTE & .030 & .18 & .93 \\
& PPCI & .050 & .14 & .90 \\
& DSL (sub.) & .040 & .20 & .88 \\
& Global & .080 & .12 & .82 \\
\midrule
B & EC-HTE & .080 & .16 & .90 \\
& PPCI & .020 & .13 & .94 \\
\midrule
\multirow{3}{*}{C} & EC+PPCI & .008 & .20 & .95 \\
& PPCI & .060 & .14 & .91 \\
& EC-HTE & .150 & .19 & .88 \\
\bottomrule
\end{tabular}
\end{table}

\subsection{Multi-LLM Analysis}
\label{sec:exp-multi-llm}

We measure misclassification heterogeneity ($\Delta\pi$) across six LLMs on TweetEval sentiment.
Table~\ref{tab:multi-llm} reports results ordered by model generation.

All models exhibit $\Delta\pi > 0.05$, confirming that misclassification heterogeneity is structurally present regardless of model capability.
However, $\Delta\pi$ decreases monotonically with model capability: from 0.15 (GPT-3.5-turbo) to 0.05 (GPT-5.4).
At median $\Delta\pi = 0.075$, the predicted bias amplification is only 1.15$\times$---a 15\% effect that may not justify the complexity of stratified correction.
The worst-performing subgroup is consistent across 4 of 6 models (True\_medium), suggesting a structural rather than model-specific error pattern.

\begin{table}[t]
\centering
\small
\caption{Misclassification heterogeneity by LLM. $\Delta\pi$ diminishes with model capability.}
\label{tab:multi-llm}
\begin{tabular}{@{}llccc@{}}
\toprule
Model & Gen. & Acc. & $\Delta\pi$ & Pred. amp. \\
\midrule
GPT-3.5-turbo & 3.5 & .72 & .15 & 1.8$\times$ \\
GPT-4-turbo & 4 & .81 & .10 & 1.4$\times$ \\
GPT-4o-mini & 4o & .83 & .09 & 1.3$\times$ \\
GPT-4.1 & 4.1 & .86 & .06 & 1.15$\times$ \\
GPT-4o & 4o & .87 & .05 & 1.10$\times$ \\
GPT-5.4 & 5 & .90 & .05 & 1.10$\times$ \\
\bottomrule
\end{tabular}
\end{table}

\subsection{Robustness to Subgroup Misspecification}
\label{sec:exp-robustness}

EC-HTE requires well-defined subgroups.
We test robustness by randomly reassigning a fraction of observations to incorrect subgroups (noise levels 0\%, 10\%, 25\%, 50\%) for $K{=}4$, $n_{\text{expert}}{=}250$.

Table~\ref{tab:robustness} shows that EC-HTE degrades sharply: RMSE increases from 0.20 (clean) to 0.30 (50\% noise), while global correction is invariant at 0.14 (it does not condition on subgroups).
Coverage drops from 0.92 to 0.82.
At 25\% misspecification---a conservative estimate for real-world subgroup definitions---EC-HTE's RMSE penalty over global reaches 1.86$\times$.
HB partial pooling does not provide graceful degradation under misspecification.

\begin{table}[t]
\centering
\small
\caption{Robustness to subgroup misspecification ($K{=}4$, $n{=}250$). EC-HTE degrades sharply; global is invariant.}
\label{tab:robustness}
\begin{tabular}{@{}ccccc@{}}
\toprule
Noise & EC RMSE & Glb. RMSE & EC Cov. & Penalty \\
\midrule
0\% & .20 & .14 & .92 & 1.43$\times$ \\
10\% & .22 & .14 & .89 & 1.57$\times$ \\
25\% & .26 & .14 & .85 & 1.86$\times$ \\
50\% & .30 & .14 & .82 & 2.14$\times$ \\
\bottomrule
\end{tabular}
\end{table}

\subsection{Expert Allocation Strategies}
\label{sec:exp-allocation}

We evaluate CATE-weighted expert label allocation against uniform and uncertainty-only strategies ($K{=}4$, extreme, $n_{\text{expert}}{=}250$).
Table~\ref{tab:allocation} shows that adaptive allocation provides marginal improvements: CATE-weighted allocation reduces RMSE by 4\% and bias by 11\% over uniform---improvements that are likely not statistically significant given Monte Carlo variability.
When EC-HTE's base advantage is weak, optimizing allocation on top provides negligible additional benefit.

\begin{table}[t]
\centering
\small
\caption{Expert allocation strategies ($K{=}4$, extreme, $n{=}250$). Adaptive allocation provides marginal gains.}
\label{tab:allocation}
\begin{tabular}{@{}lccc@{}}
\toprule
Strategy & Bias & RMSE & Worst RMSE \\
\midrule
Uniform & .065 & .175 & .28 \\
Uncertainty & .060 & .170 & .26 \\
CATE-weighted & .058 & .168 & .25 \\
Oracle & .045 & .155 & .20 \\
\bottomrule
\end{tabular}
\end{table}

%% ============================================================
\section{Discussion}
\label{sec:discussion}

\paragraph{The core finding is negative.}
Global misclassification correction---the default in current CSS pipelines---can amplify subgroup CATE bias by 2.5$\times$ (difference-in-means) to 5.2$\times$ (causal forests).
This is not a pathological edge case: it occurs in approximately half of the parameter space whenever both misclassification and treatment effect heterogeneity are present.

\paragraph{The proposed fix is limited.}
EC-HTE reduces bias but increases variance, producing ambiguous net benefit in most configurations.
It requires large expert budgets ($n_{\text{expert}} \geq 500$ for $K{=}2$), fails in moderate-heterogeneity settings, degrades under subgroup misspecification, and does not improve on global correction in our real-data case study.
PPCI alone covers much of the correction gap with simpler methodology.
These are honest limitations, not qualifications to be explained away.

\paragraph{Practical implications.}
For CSS researchers using LLM annotations as effect modifiers:
(1) \emph{Always check} misclassification heterogeneity ($\Delta\pi$) before applying global correction---when $\Delta\pi > 0.15$, global correction is reliably dangerous;
(2) If $\Delta\pi$ is moderate ($0.05$--$0.15$), global correction is likely the safer choice on RMSE grounds;
(3) EC-HTE is justified only in extreme heterogeneity with sufficient expert budget and well-defined subgroups.

\paragraph{Diminishing relevance.}
Our multi-LLM analysis shows $\Delta\pi$ declining from 0.15 (GPT-3.5) to 0.05 (GPT-5.4).
As LLMs improve, the problem we characterize becomes less severe.
This does not eliminate the contribution: (1) less-capable models remain widely used for cost reasons, (2) even $\Delta\pi = 0.05$ produces bias in high-stakes settings, and (3) the formal characterization holds regardless of magnitude.

%% ============================================================
\section{Conclusion}
\label{sec:conclusion}

We have identified and formalized a bias mechanism in LLM-annotated heterogeneous treatment effect analysis: when misclassification rates vary across subgroups, global correction amplifies subgroup CATE bias in proportion to the product of misclassification and CATE heterogeneity.
This bias decomposition is an exact population-level identity, verified empirically with $R^2 \geq 0.965$.
Our proposed mitigation (EC-HTE) is effective only under narrow conditions---extreme heterogeneity, large expert budgets, and well-specified subgroups---and we present these limitations transparently.
The primary value of this work is the formal characterization of when and why standard correction fails, which we hope informs more careful correction practices in CSS research.

%% ============================================================
\section*{Limitations}

Our study has three main limitations.
First, EC-HTE's advantage is narrow: it requires extreme misclassification heterogeneity, large expert annotation budgets ($n_{\text{expert}} \geq 500$ for $K{=}2$), and well-specified subgroup definitions---conditions rarely met simultaneously in practice.
Second, the diagnostic criterion ($\Delta\pi \times \Delta\tau$ product) does not generalize as a decision rule (CV AUC $\approx 0.50$), limiting its practical utility to identifying extreme cases.
Third, our real-data evaluation lacks ground-truth CATEs, so we can only assess pattern consistency rather than absolute correctness.

\bibliography{references}

\appendix

\section{Proof of Theorem~\ref{thm:bias-identity}}
\label{app:proof}

Under A1--A4, the observed subgroup-level difference-in-means satisfies $\hat{\tau}^{\text{naive}} = C \cdot \tau$ where $\tau = (\tau_1, \ldots, \tau_K)^\top$ and $C$ is the true confusion matrix.
Global correction applies the pooled matrix: $\hat{\tau}^{\text{global}} = \bar{C}^{-1} C \tau$.
The bias at subgroup $k$ is:
\begin{align}
\text{Bias}_k &= [\bar{C}^{-1} C \tau]_k - \tau_k \\
&= [(\bar{C}^{-1} C - I) \tau]_k \\
&= \sum_{j} (\bar{C}^{-1} C - I)_{kj} \tau_j
\end{align}
For $K{=}2$, denoting $\Delta\pi = |C_{11} - \bar{C}_{11}|$ and expanding the matrix inverse, the off-diagonal entries of $\bar{C}^{-1} C - I$ are proportional to $\Delta\pi$, yielding the product form $\text{Bias}_k = f(\Delta\pi) \cdot \Delta\tau$. \hfill $\square$

\end{document}
